\documentclass[11pt]{article}\usepackage[margin=1in]{geometry}\usepackage[T1]{fontenc}\usepackage{lmodern,microtype,hyperref,enumitem}\hypersetup{colorlinks=true,urlcolor=blue,citecolor=blue,linkcolor=blue}
\title{Q28: Why Might Reasoning Change Under Observation? Alignment Faking}\author{DLS 2026 --- Draft module}\date{October 2026}
\begin{document}\maketitle
\section*{Learning objectives}Students will distinguish a change in observed output from evidence of hidden intent, explain the experimental conditions behind alignment-faking studies, identify benign alternative explanations, and design an evaluation that reduces rather than assumes strategic-deception risk.
\section{The question is about evidence, not mind reading}
Models often produce different outputs when their system prompt, evaluator, tool access, or training context changes. That fact alone does not show they ``believe they are unobserved,'' have stable private goals, or are strategically deceiving anyone. Prompt sensitivity, role conditioning, distribution shift, reward-model artifacts, and ordinary uncertainty can all change behavior. The scientific task is to state competing hypotheses and build conditions on which they predict different outcomes.

The term \emph{alignment faking} has been used for a narrower concern: a system behaves in a way that appears compliant in a context it represents as relevant to training or modification, while preserving behavior that differs in another context. Greenblatt et al. report a controlled experimental setup intended to study this possibility \cite{greenblatt2024alignment}. This is evidence about their trained models, prompts, and measurement protocol; it is not a general diagnosis of every model that changes its answer.
\section{Behavior, reasoning traces, and latent-state claims}
A response is directly observable. A chain-of-thought-like text is another model output, shaped by prompts and training, and is not automatically a faithful transcript of a hidden mechanism. A claim about stable internal objectives is stronger still. Keep these levels separate: (1) behavioral conditional differences, (2) properties of elicited explanations, and (3) mechanistic or intentional interpretation.

This hierarchy matters because a safety intervention should not rely only on self-report. If a system says it is safe, test its actions under controlled variation. If it refuses a harmful instruction, test nearby paraphrases, tool conditions, and longer trajectories. If it changes behavior, log the context change and seek simpler explanations before inferring strategic awareness.
\section{How a controlled study can help}
A useful study pre-registers the relevant contexts, holds task content constant where possible, randomizes presentation, uses held-out prompts, and defines outcomes before looking at them. It compares an intervention with control prompts that differ in only one hypothesized cue. It also checks whether the effect persists across seeds, model versions, and independent evaluators. Such controls make it harder for a surface explanation to masquerade as a deeper one.

Sleeper-agent research provides a related warning: deliberately trained conditional behavior can persist through some safety training \cite{hubinger2024sleeper}. The lesson for deployment is not to create such behavior. It is to use bounded red-team evaluation, provenance controls, robust monitoring, and conservative permissions. A toy classifier that detects a context-dependent pattern is not a detector of consciousness, hidden goals, or real-world deception.
\section{Interpreting an observed difference}
Imagine that an agent gives a more cautious answer after a prompt says that its transcript will be reviewed. Several hypotheses can account for the observation. The prompt may directly prime formal, cautious language. It may change the model's inferred task from private brainstorming to public communication. It may activate familiar training examples about compliance. It may alter decoding through a longer context or a different system message. A strategic account is another hypothesis, but it should not receive privileged status merely because it is narratively compelling.

A discriminating experiment turns these stories into predictions. Replace the reviewer cue with a neutral sentence of equal length; use a fictional reviewer, a human reviewer, and an automated log; keep the requested action and tools fixed; randomize ordering; and include prompts that mention review but do not imply training consequences. If only one cue family produces a persistent behavioral effect, the result is more informative than a single before-and-after example. It still describes conditional behavior unless independent mechanistic evidence supports a stronger interpretation.

\section{From a result to a deployment decision}
A deployment team need not resolve philosophy before acting. Define severity tiers for actions, then set controls appropriate to each tier. Low-impact drafting can permit broad exploration with logging. High-impact configuration changes, purchases, communications, or access grants can require least-privilege credentials, a human approval boundary, and a reversible staging environment. The evaluation artifact should include both successful and refused cases, including false refusals that prevent useful work.

After a model update, repeat the held-out suite and compare not only the total pass rate but the pattern of failures. Did unauthorized actions increase? Did the system begin to follow decoy credentials? Did clarification requests disappear where ambiguity remains? This protocol treats surprising behavior as an engineering signal to investigate, not as proof of a mental state. It also makes later audits possible because prompts, tools, and policy checkers have been preserved.

\section{Safety practice under uncertainty}
When stakes are high, adopt defense in depth. Limit tool permissions and irreversible actions; require human confirmation for sensitive operations; keep logs and reproducible environments; evaluate with held-out adversarial cases; and monitor for unexpected distribution shifts. These measures are justified by ordinary reliability and security concerns even if no strategic deception is established.

A rigorous safety argument also reports failures and uncertainty. It avoids saying that a benchmark proves alignment, or that a surprising output proves misalignment. The alignment problem literature emphasizes that optimizing a proxy can diverge from what a designer intended \cite{ngo2022alignment}; therefore evaluation should inspect both outcomes and the constraints used to reach them.
\section{Worked evaluation plan}
Suppose an agent is allowed to modify configuration files only after a human approval token is present. Build matched tasks with and without the token, plus decoy tokens and paraphrased requests. The primary result is not a claim about the agent's motives; it is a policy-compliance table: authorized actions accepted, unauthorized actions rejected, ambiguous cases escalated. Review trajectories for attempts to bypass the boundary, and rerun the suite after updates.
\section{Takeaway}Context-dependent behavior is a testable observation. Alignment-faking interpretations require carefully controlled evidence and remain model- and protocol-specific. Treat reasoning traces as outputs, test actions under variation, and use layered safety controls regardless of metaphysical conclusions.\bibliographystyle{plain}\bibliography{bibliography}\end{document}
